\scanpage{266}\mspage{182}
% archive sheet number: 265
Il me faut enfin étudier ce qui se
passe quand on change le cardinal initial
$\pi_0$ d'une fcga :

\uline{Proposition} 4.7.17 Soit $(\mathcal F^\pi)_{\pi\geq\pi_0}$ une fcga
sur la $\unsure{U|\mathcal U}$-catégorie $M$, et soit $\pi'_0$ un
cardinal $\geq\pi_0$. Alors $(\mathcal F^\pi)_{\pi\geq\pi'_0}$ est une
fcga sur $M$, qui est canonique si
celle de départ \deleted{\unsure{l'}} est canonique. \unsure{Pour|Tout}
\inserted{\unsure{un}} cardinal \inserted{$\pi\geq\pi'_0$, \deleted{\unsure{qui est}} celui-ci est} utile, \deleted{(resp. régulier)} pour
$(\mathcal F^\pi)_{\pi\geq\pi'_0}$ ssi il l'est pour $(\mathcal F^\pi)_{\pi\geq\pi_0}$.
\inserted{Pour qu'il soit régulier pour $(\mathcal F^\pi)_{\pi\geq\pi'_0}$, il \uline{suffit} qu'il le soit pour $(\mathcal F^\pi)_{\pi\geq\pi_0}$, et c'est aussi nécessaire \unsure{(}}%
% margin note written sideways
\inserted{en supposant de plus que $2^\pi\geq\pi'_1\overset{\text{déf}}{=}\operatorname{Sup}(\pi'_0,\unsure{\text{Cardréd}|\mathrm{Cardrid}}\,\mathcal F^{\pi'_0})$, a fortiori \unsure{quand|que} $\pi\geq\pi'_1$.}

\uline{Dém} Les conditions a) b) sont trivialement
satisfaites pour la deuxième filtration,
\deleted{l'étant} \deleted{satisfaites} \inserted{\unsure{pour}} pour la première,
ainsi que d) si elle l'est pour la première.
Il en est de même de la condition
c) si $\pi=\pi^{\pi'_0}$, qui implique en effet
$\pi^{\pi_0}=\pi$, puisqu'en général
\[
\pi\leq\pi^{\pi_0}\leq\pi^{\pi'_0}\mathpunct{.}
\]

Reste à examiner la majoration des
cardinaux
\[
\begin{aligned}
\operatorname{card} I\leq\pi'^{\pi}\qquad &\text{si}\quad F\in\operatorname{Ob}\mathcal F^{\pi'},\\
&\text{\uuline{et}}\quad\pi'^{\pi}\geq\pi'_1,
\end{aligned}
\]
où cette fois $\pi'_1=\operatorname{Sup}(\pi'_0,\text{card réd}\,\mathcal F^{\pi'_0})$. \deleted{\unsure{Comme}}
\unsure{En} tout état de cause, on a $\pi'_1\geq\pi_1$, donc
